\documentclass[utf8]{frontiers_suppmat}
\usepackage{url,hyperref,lineno,microtype}
\usepackage[onehalfspacing]{setspace}
\usepackage{booktabs}

\begin{document}
\onecolumn
\firstpage{1}

\title[Supplementary Material]{{\helveticaitalic{Supplementary Material: Full Analysis of the SIMPLEX Case Study}}}

\maketitle

**SIMPLEX: Composition-Space Deep Learning for Hydrogel Adhesion with
Out-of-Distribution Extrapolation to Model-Discovered High-Performance
Formulations**

Tongfei Shen

---

\subsection{S1. Dataset description}

\subsubsection{S1.1 Source and provenance}

The study uses the publicly released dataset accompanying Liao et al.,
*Nature* 644, 89–95 (2025), DOI 10.1038/s41586-025-09269-4
(repository: \texttt{sheng-hu/hydrogels}, MIT licence). The dataset records
underwater glass adhesion strength (kPa) for 341 bio-inspired hydrogel
formulations, each described by the molar fractions of six functional
monomers:

\begin{tabular}{llll}
\toprule
# & Monomer & Functional class & Abbreviation \\
\midrule
1 & 2-Hydroxyethyl acrylate & Nucleophilic & HEA \\
2 & Butyl acrylate & Hydrophobic & BA \\
3 & 2-Carboxyethyl acrylate & Acidic & CBEA \\
4 & (3-Acrylamidopropyl)trimethylammonium chloride & Cationic & ATAC \\
5 & 2-Phenoxyethyl acrylate & Aromatic & PEA \\
6 & Acrylamide & Amide & AAm \\
\bottomrule
\end{tabular}


The six fractions sum to 1 (composition simplex). Target: \texttt{Glass (kPa)_max}.

\subsubsection{S1.2 Internal / external split (model-guided extrapolation)}

\begin{itemize}
\item \textbf{Internal (training region):} the 180 round-1 bio-inspired formulations
\end{itemize}
  (\texttt{df_180.csv}). Adhesion mean 46.9 kPa, max 146.6 kPa.
\begin{itemize}
\item \textbf{External (extrapolation region):} the 161 formulations added in later
\end{itemize}
  iterations of the Nature study's sequential model-based optimisation (SMBO)
  loop (\texttt{df_341.csv} minus \texttt{df_180.csv}). Adhesion mean 154.2 kPa, max
  353.3 kPa; 45% of external samples exceed the training maximum.

This is a *model-guided migration to a high-performance composition region*
(target-value extrapolation), not an independent-laboratory cohort. The
selection bias and range restriction inherent to the SMBO protocol are
acknowledged; range-restricted Spearman rho is therefore a conservative
estimate of ranking ability.

\subsubsection{S1.3 Features}

Two modalities, 21 features total:

1. \textbf{monomer_fractions (6):} raw molar fractions.
2. \textbf{pairwise_synergy (15):} all products x_i·x_j (i<j), explicitly encoding
   monomer-pair interactions.

Scaling/imputation were fitted inside each CV training fold only
(no leakage).

---

\subsection{S2. Model configuration (final, best_config.json)}

```
d_model=64, n_blocks=2, n_heads=4, dropout=0.2
use_transformer=False (ResBlock base), use_attention=True
use_film=False, use_modality_gate=False
use_mixup=True (alpha=0.4), use_swa=True
use_sam=False, use_ema=False
use_contrastive=False, use_pretrain_recon=False, use_mfm=False
use_uncertainty_weighting=False (single target)
use_domain_constraint=True (constraint_w=0.1)
y_transform=standard, scaler=standard
max_epochs=150, patience=30, batch_size=32
lr=3e-3, weight_decay=1e-3
```

Design notes:
\begin{itemize}
\item \textbf{ResBlock base, not Transformer:} on 180 samples the Transformer
\end{itemize}
  diverged (internal R2 ≈ -6); the residual network with interaction
  attention is the stable configuration (internal R2 ≈ 0.71).
\begin{itemize}
\item \textbf{Mixup is the key regulariser:} removing it drops internal R2 from
\end{itemize}
  0.713 to 0.646.
\begin{itemize}
\item \textbf{Self-supervised pre-training (MFM / contrastive) was detrimental}
\end{itemize}
  (internal R2 ≈ -2.7) and was removed (ablation-gated).

---

\subsection{S3. Full ablation table (internal R2, 2 seeds x 3 folds)}

\begin{tabular}{ll}
\toprule
Variant & delta R2 (full − variant) \\
\midrule
w/o multimodal fusion & +0.093 (largest) \\
w/o Mixup & +0.067 \\
fusion = gated & +0.062 \\
w/o residual blocks & +0.062 \\
fusion = film & +0.031 \\
w/o task-specific gating & +0.022 \\
w/o attention & +0.019 \\
w/o domain constraint & +0.000 (neutral) \\
w/o SWA & +0.008 \\
... (full 20-component ablation in results/ablation/ablation.csv) \\
\bottomrule
\end{tabular}


Positive delta = removal degrades performance (component helps). 16
components that did not pay for themselves were removed from the final
model; the remaining switches are reported transparently.

---

\subsection{S4. External ranking statistics}

Sample-level bootstrap (n=161, 2000 resamples):

\begin{tabular}{lll}
\toprule
Model & Spearman rho & 95% CI \\
\midrule
SIMPLEX & 0.501 & [0.369, 0.619] \\
ElasticNet & 0.494 & [0.36, 0.62] \\
Ridge & 0.486 & [0.35, 0.61] \\
SVR-RBF & 0.379 & [0.24, 0.51] \\
MLP & 0.315 & [0.18, 0.45] \\
RandomForest & 0.211 & [0.06, 0.36] \\
\bottomrule
\end{tabular}


Paired bootstrap (SIMPLEX − baseline):
\begin{itemize}
\item vs RandomForest: delta rho = +0.183, 95% CI [+0.066, +0.311],
\end{itemize}
  P(diff>0)=0.999 $\rightarrow$ \textbf{significant}
\begin{itemize}
\item vs ElasticNet: delta rho = −0.028, 95% CI [−0.157, +0.097],
\end{itemize}
  P(diff>0)=0.345 $\rightarrow$ \textbf{statistical tie} (reported honestly)

Top-k screening precision (fraction of predicted top-k that are true top-k):

\begin{tabular}{llll}
\toprule
k & SIMPLEX & RF & ElasticNet \\
\midrule
20 & 0.25 & 0.10 & 0.05 \\
30 & 0.37 & 0.07 & 0.17 \\
\bottomrule
\end{tabular}


Prediction-range behaviour on the external set (kPa):
\begin{itemize}
\item SIMPLEX [16, 109]; RandomForest [26, 109] (compressed — tree models are
\end{itemize}
  locked near the training maximum ~147); ElasticNet [−82, 232] (over-extrapolating,
  non-physical negative predictions); true range [3, 353].
\begin{itemize}
\item SIMPLEX is the only model whose external predictions stay physically
\end{itemize}
  plausible (non-negative, in-range) while retaining ranking signal.

---

\subsection{S5. Multi-target extensibility (preliminary)}

The original dataset also records rheological targets for the 180 internal
formulations: storage-related modulus (kPa) and loss tangent tan-delta.
\texttt{corr(Glass, Modulus) = −0.036} — independent information not used by the
Nature study's single-target models. As a preliminary validation of
multi-target extensibility, a single-task SIMPLEX trained on Modulus
achieves internal CV R2 = 0.396 vs RandomForest 0.392 (3 seeds), i.e. parity
with the strongest baseline on a second, independent target. Joint
multi-target training is left for future work (see Limitations).

---

\subsection{S6. Interpretation highlights}

Cross-validated permutation importance (mean over folds/seeds):

\begin{tabular}{llll}
\toprule
Rank & Feature & Importance & SE \\
\midrule
1 & Cationic-ATAC & 0.222 & 0.028 \\
2 & pair (ATAC$\times$PEA) & 0.099 & 0.017 \\
3 & pair (BA$\times$CBEA) & 0.072 & 0.010 \\
4 & Hydrophobic-BA & 0.055 & 0.013 \\
5 & pair (BA$\times$PEA) & 0.050 & 0.020 \\
\bottomrule
\end{tabular}


The cation monomer and pairwise interaction terms dominate — consistent
with the electrostatic/hydrophobic synergy known to drive underwater
adhesion, and with the SMBO-discovered high-performance region enriched in
the hydrophobic–aromatic (BA–PEA) combination.

---

\subsection{S7. Reproducibility}

\begin{itemize}
\item Environment: Python 3.11 (conda env \texttt{HydroGelNet}), PyTorch 2.14.0+cu130,
\end{itemize}
  NVIDIA RTX 5080 16 GB.
\begin{itemize}
\item Fixed seeds: [42, 2024, 7, 1337, 20260731]; all preprocessing inside folds.
\item Full pipeline: `download $\rightarrow$ build $\rightarrow$ qc $\rightarrow$ tune $\rightarrow$ train $\rightarrow$ baselines $\rightarrow$ stats
\end{itemize}
  $\rightarrow$ gate $\rightarrow$ interpret $\rightarrow$ figures $\rightarrow$ tables $\rightarrow$ paper` (run_all.py).
\begin{itemize}
\item PERF-GATE verdict: \textbf{PASS} (G1 internal tie, G2 no significant
\end{itemize}
  disadvantage, G3 seed stability, G4 external Spearman 0.501 vs best
  baseline 0.494, G5 informative ablation).

---

\subsection{S8. Data & code availability}

\begin{itemize}
\item Data: \texttt{sheng-hu/hydrogels} (MIT), Nature 2025 supplement.
\item Code: https://github.com/<owner>/HydroGelNet (to be released).
\item Reproducible artefacts: \texttt{data/}, \texttt{results/}, \texttt{figures/}, \texttt{tables/},
\end{itemize}
  \texttt{paper/} in the project repository.


\end{document}
